\documentclass[12pt]{article}

\usepackage[top=1in,bottom=1in,left=1.25in,right=1.25in]{geometry}
\usepackage{amsmath,amssymb,amsthm}
\usepackage{graphicx}
\usepackage{booktabs}
\usepackage{algorithm}
\usepackage{algpseudocode}
\usepackage{hyperref}
\usepackage{xcolor}
\usepackage{listings}
\usepackage{caption}
\usepackage{subcaption}
\usepackage{enumitem}
\usepackage{url}
\usepackage{float}

\hypersetup{
    colorlinks=true,
    linkcolor=blue!70!black,
    citecolor=green!60!black,
    urlcolor=blue!70!black,
}

\lstset{
    basicstyle=\ttfamily\small,
    breaklines=true,
    frame=single,
    backgroundcolor=\color{gray!10},
}

\title{\textbf{Closed-Loop Streaming Synthesis}\\[2pt]
\textbf{Bounded-Drift, Arbitrary-Length Audio--Video Generation}\\
\textbf{by Treating Chunk Hand-off as Feedback Control}}

\author{Aleksander Majda\\
}

\date{August 2026}

\begin{document}

\maketitle

\begin{abstract}
Generating arbitrarily long videos with diffusion transformers requires splitting the
temporal dimension into overlapping chunks.  Naive open-loop chunk-streaming
accumulates exposure-bias drift: each chunk is conditioned on the previously
generated latent, compounding distributional errors and causing scene collapse or
grain amplification after a few hundred frames.  We present
\textbf{Closed-Loop Streaming Synthesis (CLSS)}, a suite of closed-loop corrections
that close the feedback loop between successive chunks without modifying the
transformer weights: calibrated context re-noising, per-channel EMA-tracked
AdaIN, frequency-band soft shrinkage, a dynamic anchor bank, and a two-band
spatial detail anchor.  CLSS is implemented as a ComfyUI custom-node pipeline
targeting LTX-2.3 22B (a 44\,GB audio-video diffusion transformer) on consumer
16\,GB VRAM hardware via GGUF 4-bit quantisation and CPU block-streaming, driven
by a multi-modal guider with split per-modality classifier-free guidance and
spatiotemporal guidance via a value-passthrough self-attention skip, with all
streaming noise seeded for bit-reproducible runs.  A two-stage pipeline
generates a streaming Stage-1 draft and refines it at $2\times$ spatial
resolution with a distilled schedule.

We report the production system at its operating point: single-scene,
$69.3$-second clips streamed as four chunks of 52 latent frames ($\approx17.3$\,s
windows, inside the model's $\approx20$\,s positional-encoding wall), text-only
(T2V) and image-to-video (I2V) variants, with the final outputs verified
perceptually.  Two system-level findings are validated.  First, the
production chunk length carries a sigma-schedule conflict: the connected
scheduler derives its shift from the chunk template's token count, and at 52-lf
chunks the video requires the resulting $4.68$ shift while the audio requires
the $1.844$ shift --- no shared schedule serves both modalities, so the
pipeline splits them inside the single joint pass (audio shift scaled by
$0.394$, shape-only, full-noise start).  Second, Stage-2 audio refinement is
joined at a sigma cutoff: re-inventing audio from the Stage-2 noise level
destroys it because melody and timbre are not text-specified, while a join at
$\sigma=0.55$ preserves $45\%$ of the Stage-1 audio and adds the low-$\sigma$
descent the audio would otherwise never receive.  An open/closed-loop ablation
at the production geometry confirms the corrections are what bound drift:
without them video std rises $+0.118$ over five chunks, audio RMS $+34\%$,
and audio high-frequency energy drifts to $1.60\times$ its chunk-1 reference;
with them video std holds to $+0.001$ and object identity holds
$0.980$--$0.995$ (per-cell minimum $0.79$).  In the verified production runs,
Stage-1 video boundary similarity holds $0.93$--$0.99$, the decoded seam-grid
object identity holds $0.988$--$0.998$ mean cell cosine, video latent std
drifts $<0.005$, the audio RMS anchor holds loudness to its chunk-1 reference
($0.436\to0.504$ upward-only), and the Stage-2 joined refinement holds
audio fidelity to Stage-1 at $0.66$--$0.74$.  Coarse scene layout is
monitored by a sign-free power-map metric, and a per-chunk telemetry suite
(settings dumps, per-frame drift-event tables, phase-lock detection)
localises every reported failure in wall-clock time and latent-frame
position.
\end{abstract}

%--------------------------------------------------------------------------
\section{Introduction}
\label{sec:intro}
%--------------------------------------------------------------------------

Modern video diffusion transformers operate on patchified latent representations
where temporal resolution is down-scaled $8\times$ relative to pixel frames.
Fitting a 30-second clip into a single forward pass requires very large GPUs.
The standard remedy is \emph{temporal chunking} via a Streaming Latent Buffer
(SLB): generate the video in short overlapping segments, conditioning each on
the previous denoised output.  Section~\ref{sec:background} reviews this
mechanism and the exposure-bias problem it introduces.

CLSS closes the feedback loop between chunks by applying five lightweight
post-denoising corrections that together bound the closed-loop gain below one,
preventing drift from accumulating without modifying the transformer weights.
The corrections are applied purely at the boundary between chunks and add
negligible computational overhead.

\paragraph{Contributions.}
\begin{enumerate}[leftmargin=*,noitemsep]
\item Five closed-loop correction modules for temporal chunk streaming:
      calibrated context re-noising, EMA-tracked per-channel AdaIN,
      frequency-band soft shrinkage, a dynamic anchor bank, and a two-band
      spatial detail anchor that closes the progressive-softening channel the
      one-sided shrinkage leaves open
      (Sections~\ref{sec:renoise}--\ref{sec:detail}).
\item A linear stability analysis of the feedback loop in latent space with an
      empirical per-chunk gain measurement protocol.  The measurement
      resolves an apparent paradox: the one-step boundary sensitivity
      ($g_{\rm SLB}\approx15$--$87$) sits an order of magnitude above the
      nominal stability boundary, yet the closed loop is observably stable ---
      the contraction back toward the data manifold is what bounds drift
      (Sections~\ref{sec:stability}, \ref{sec:discussion}).
\item A per-modality sigma-schedule split inside a single joint diffusion pass,
      resolving the measured conflict at the production chunk length where the
      video and the audio require different schedule shifts
      (Sections~\ref{sec:sigmasplit}, \ref{sec:exp_split}).
\item A Stage-2 joined audio refinement: the audio re-denoises on the same
      low-$\sigma$ descent as the video, joined at a sigma cutoff that keeps
      the unprompted content intact (Sections~\ref{sec:s2audiorefine},
      \ref{sec:exp_join}).
\item A memory-efficient two-stage pipeline for LTX-2.3 22B on 16\,GB VRAM
      using GGUF 4-bit quantisation and CPU block-streaming
      (Section~\ref{sec:impl}).
\item A multi-modal guider for ComfyUI (split per-modality CFG, STG via
      value-passthrough self-attention skip) that correctly handles ComfyUI's
      packed audio-video latent representation (Section~\ref{sec:guider}).
\item ComfyUI custom nodes implementing the full pipeline with seeded,
      bit-reproducible streaming noise and per-chunk telemetry: unconditional
      settings dumps, per-frame origin/layout drift event tables, a phase-lock
      repetition detector, and a spatially-resolved seam identity metric
      (Sections~\ref{sec:nodes}, \ref{sec:telemetry}).
\item Production-configuration results verified perceptually: matched
      $69.3$-second T2V and I2V generations at 52-lf chunks with the
      per-modality split and joined Stage-2 refinement, an open/closed-loop
      ablation at the same geometry, a sign-free coarse power-map layout
      metric, and an event-triggered seam-pin safeguard
      (Sections~\ref{sec:experiments}, \ref{sec:events}).
\end{enumerate}

%--------------------------------------------------------------------------
\section{Background}
\label{sec:background}
%--------------------------------------------------------------------------

\subsection{LTX-2.3 Architecture}
LTX-2.3 is a 22-billion-parameter diffusion transformer for joint audio-video
synthesis.  It operates on a latent space with an $8\times$ temporal downscale
(video VAE, causal) and a $160\times4\times$ audio downscale (audio VAE,
16\,kHz, hop length 160, compression factor 4), yielding a video latent rate
of $25/8 = 3.125$ lf/s and an audio latent rate of $16000/(160\times4)=25.0$
lf/s.  The model accepts interleaved video and audio tokens together with
Gemma-3 text embeddings and is trained with flow-matching on the
rectified-flow schedule.

\paragraph{The positional-encoding wall.}
Audio-video RoPE positions are measured in \emph{seconds}, computed per forward
pass with $\text{max\_pos}=20$ for both modalities.  Any single window longer
than $\approx20$\,s runs off the positional grid: measured audio glitches at
$\approx20.5$\,s (a transient at latent frame 512--513, at the same position
across seeds and resolutions) and video goes near-static.  Every streaming
chunk is therefore kept below the wall; the production geometry uses
52-lf windows ($\approx17.3$\,s new content plus overlap $\approx2.7$\,s
$=20.0$\,s total), which is the largest window that fits.

\subsection{Streaming Latent Buffer (SLB)}
The SLB provides temporal continuity in chunk streaming.  When generating
chunk $N$, the last $\text{overlap\_lf}$ latent frames from chunk $N-1$ are
prepended to the new chunk's latent, and the denoiser is conditioned on them
at noise level $\tau_c$:
\begin{equation}
\tilde{L}_{\text{overlap}} = (1-\tau_c)\,\hat{L}_{\text{overlap}} + \tau_c\,\varepsilon,
\quad \varepsilon \sim \mathcal{N}(0, I)
\label{eq:renoise}
\end{equation}
Setting $\tau_c=0$ yields maximal continuity but maximal drift risk;
$\tau_c=1$ decouples the chunks entirely.  In LTX's framework this is
implemented as \texttt{VideoConditionByLatentIndex} with
$\text{strength}=1-\tau_c$.

\paragraph{Implementation note: masked timestep, not a one-shot draw.}
Eq.~\eqref{eq:renoise} describes the target distribution of the overlap
context, but the sampler does not literally draw $\varepsilon$ once and hand
the model a truncated schedule from $\tau_c$.  Tracing the actual mechanism
shows that the per-token timestep supplied to the transformer at every
sampling step $t$ is $\tau_{\text{eff}} = (1-\text{strength})\cdot\sigma_t
= \tau_c\,\sigma_t$ for overlap tokens, versus $\sigma_t$ for the rest of the
sequence --- a \emph{persistent} per-step rescaling of the global noise
schedule, applied at every step, not a single noise-injection event.  For a
\emph{linear} (rectified-flow) interpolation path,
$x_t = (1-t)x_0 + t\varepsilon$, the two descriptions coincide exactly:
scaling every step's effective time by a constant factor $\tau_c$ is
algebraically identical to drawing $x_{\tau_c}$ once via
Eq.~\eqref{eq:renoise} and continuing the \emph{same} linear ODE from
$\tau_c$ down to $0$.  LTX-2.3 is trained with a rectified-flow objective, so
the equivalence holds for this model; it would not hold in general for a
non-linear noise schedule.

\subsection{Exposure Bias in Streaming Diffusion}
Let $\hat{L}_N$ be the denoised latent of chunk $N$ and
$f_\theta(\cdot\mid\hat{L}_{N-1})$ the denoising operator conditioned on the
previous chunk.  In open-loop streaming:
\begin{equation}
\hat{L}_N = f_\theta(\hat{L}_{N-1} + \varepsilon_N)
\end{equation}
The transformer was trained on latents drawn from the data distribution, but at
inference it receives denoised latents from prior chunks, which lie
off-manifold relative to its training distribution.  Any distributional shift
$\delta_N = \hat{L}_N - L_N^*$ (where $L_N^*$ is the ``on-manifold'' latent
for the same content) is in general amplified by the open-loop gain of the
denoising process.  Formally, if $g = \|\partial f_\theta / \partial L\|$
evaluated at $L_{N-1}^*$, then $\|\delta_{N+1}\| \approx g\|\delta_N\|$.
When $g > 1$ this gives $\|\delta_N\| \sim g^N \to \infty$.

In practice $g$ is not a global constant: it depends on the noise level
schedule, step count, guidance strength, and the spectral content of the
latent.  Section~\ref{sec:stability} describes how $g$ is estimated
empirically per chunk via a perturbation measurement.

%--------------------------------------------------------------------------
\section{Closed-Loop Streaming Synthesis}
\label{sec:clss}
%--------------------------------------------------------------------------

CLSS adds five closed-loop correction modules that together bound the error
accumulation by ensuring $\rho_{\text{closed}} = g \cdot (1-\beta)\cdot(1-\tau_c) < 1$.

\subsection{Calibrated Context Re-Noising}
\label{sec:renoise}

The overlap latent fed as context to chunk $N$ is re-noised to level $\tau_c$
via Eq.~\eqref{eq:renoise}.  This forces the denoiser to actively re-project
the boundary context onto the data manifold rather than accepting it verbatim,
partially correcting accumulated drift in the SLB.

In practice $\tau_c=0.05$ provides a strong temporal constraint
(boundary cosine similarity $>0.97$) while allowing enough distributional
repair to prevent drift accumulation.  Larger values ($\tau_c=0.20$, as
suggested by some prior work) loosen the constraint too much for 22B GGUF
streaming, causing visible motion discontinuities.

\subsection{EMA-Tracked Per-Channel AdaIN}
\label{sec:adain}

A slow exponential moving average (EMA) with rate $\lambda \approx 0.10$
tracks per-channel mean $\mu_c$ and standard deviation $\sigma_c$ across chunks:
\begin{align}
\mu_c^{(N)} &= (1-\lambda)\,\mu_c^{(N-1)} + \lambda\,\hat\mu_c^{(N)} \\
\sigma_c^{(N)} &= (1-\lambda)\,\sigma_c^{(N-1)} + \lambda\,\hat\sigma_c^{(N)}
\end{align}
where $\hat\mu_c^{(N)}, \hat\sigma_c^{(N)}$ are the per-channel statistics of
chunk $N$'s new frames.  After generation, the overlap region is dropped and
\emph{only the new frames} $L_N^{\text{new}}$ (Algorithm~\ref{alg:clss},
line~5) are blended toward the EMA reference via AdaIN with blend factor
$\beta$:
\begin{equation}
\tilde{L}_N^{\text{new}} = (1-\beta)\,L_N^{\text{new}} + \beta\;\text{AdaIN}(L_N^{\text{new}} \to \mu^{(N-1)}, \sigma^{(N-1)})
\end{equation}
The overlap region itself is never re-touched by AdaIN --- it was already
corrected as \emph{previous} chunk's $L_{N-1}^{\text{new}}$ before being pushed
to the SLB.  This suppresses fast per-chunk statistical drift while allowing
slow intentional scene evolution (lighting changes, colour grading) to pass
through.

Two safety caps prevent over-correction:
\begin{itemize}[noitemsep]
\item \textbf{Sigma max drift}: the EMA $\sigma_c$ is clamped to at most
  $(1+\Delta)\times\sigma_c^{(0)}$ (default $\Delta=0.05$), preventing AdaIN
  from amplifying late-chunk residual noise.
\item \textbf{AdaIN max amplification}: per-channel upward scaling is capped
  at $1.2\times$ the current chunk's std (default), preventing AdaIN from
  boosting noise in channels that have become quiet.
\end{itemize}

\subsection{Frequency-Band Soft Shrinkage}
\label{sec:shrink}

The latent is decomposed into $B=4$ frequency bands via 3-D FFT along
the (frame, height, width) axes.  Each band $b$ receives a continuous gain:
\begin{equation}
g_b^{(c)} = \min\!\left(1,\; \left(\frac{E_{\text{ref},b}^{(c)}}{E_{b}^{(c)}}\right)^{\!\gamma_b}\right) \in (0,1]
\label{eq:gain}
\end{equation}
where $E_b^{(c)}$ is the mean squared amplitude of band $b$ in channel $c$,
$E_{\text{ref},b}^{(c)}$ is the corresponding EMA reference, and
$\gamma_b$ is the per-band shrinkage exponent.

Only bands with excess energy relative to their EMA are attenuated.  Bands with
$\gamma_b = 0$ (the DC band) are never attenuated.  Default exponents:
$\boldsymbol\gamma = (0.0,\,0.05,\,0.2,\,0.3)$, which are conservative enough
to preserve scene structure while suppressing temporal flicker.

The EMA is updated with the \emph{corrected} (post-gain) energies, not the raw
ones, so the reference tracks intentional scene content rather than accumulated
error.

\subsection{Dynamic Anchor Bank}
\label{sec:anchor}

A small bank of $M \le 8$ keyframe anchors provides long-range visual identity
conditioning.  The bank is seeded from the \emph{last} frame of chunk 1 (not
the first) and updated when two consecutive chunks trigger a scene-change
signal (cosine similarity below threshold $\rho_a=0.78$ for $P=2$ consecutive
chunks).  A \emph{forced insertion} every $F_a$ chunks prevents bank
stagnation on visually stable sequences where the similarity threshold is
never triggered.  The ComfyUI node auto-derives
$F_a = \max(2, \min(5, \lceil\text{num\_chunks}/4\rceil))$ ($F_a=2$ for the
four-chunk production runs).

Anchors are retrieved via cosine similarity.  In the current pipeline the bank
serves as the identity-telemetry reference (each chunk is compared against its
nearest banked anchor per chunk); anchor \emph{injection} into the
conditioning was tested in two wirings and both produced visible artifacts
(appended keyframe tokens corrupt the joint audio-video conditioning; an
in-place replacement at the first new frame morphs the content back toward
earlier scene content, because the redundancy filter guarantees the retrieved
anchor differs from the content being continued).  The bank therefore measures
identity drift without correcting it; a safe anchor-conditioning mechanism is
an open design problem (Section~\ref{sec:events}).

\subsection{Two-Band Spatial Detail Anchoring}
\label{sec:detail}

The shrinkage of Section~\ref{sec:shrink} is deliberately one-sided
(Eq.~\eqref{eq:gain} only ever attenuates), and this asymmetry has a measured
long-run consequence: low-band energy inflates while high bands are only ever
shrunk, so the \emph{high-frequency share} of latent energy falls monotonically
--- $+19.6\%$ band-0 inflation over seven chunks in one measured run, with the
uncapped band EMA itself drifting upward and legitimising the excess.  Nothing
in the shrinkage loop can \emph{restore} band energy, and the perceptual result
is progressive detail loss.

The detail anchor closes exactly this channel.  Each frame is split into a
spatial low band (a $3\times3$ average pool over the latent grid) and its
high-band residual, and the mean squared energy of each band is equalised
toward a \emph{fixed} scene-first reference (captured at chunk~1, re-baselined
on scene change) with symmetric per-chunk gains:
\begin{equation}
g_{\rm lo} = \Big(\tfrac{E_{\rm lo}^{\rm ref}}{E_{\rm lo}}\Big)^{1/2}
\;\text{clamped to } [0.90, 1.10], \qquad
g_{\rm hi} = \Big(\tfrac{E_{\rm hi}^{\rm ref}}{E_{\rm hi}}\Big)^{1/2}
\;\text{clamped to } [0.90, 1.12]
\label{eq:detail}
\end{equation}
Unlike the shrinkage EMA, the reference here does not track the output ---
drift cannot legitimise itself --- and unlike the shrinkage gains, these are
symmetric: lost high-band energy is restored, not merely excess removed.  The
per-chunk clamp bounds the correction so a genuine scene change cannot be
fought harder than $\pm10$--$12\%$ per chunk, and a $\pm0.5\%$ dead-band makes
the module a strict no-op when energies already match.  The same mechanism runs
per-window in Stage~2, referenced there to the upscaled Stage-1 slice of the
same window (Section~\ref{sec:twostage}).  In the verified production runs it
holds the Stage-1 high-frequency share flat end-to-end ($0.527\to0.527$ I2V,
$0.616\to0.616$ T2V; Section~\ref{sec:experiments}).

\subsection{Stability Analysis}
\label{sec:stability}

The AdaIN blend (Section~\ref{sec:adain}) and context re-noising
(Section~\ref{sec:renoise}) together reduce the magnitude of any per-chunk
latent perturbation by a factor $(1-\beta)(1-\tau_c)$ before it reaches the
next chunk's denoiser.  This gives the \emph{loop gain}:
\begin{equation}
\rho_{\text{loop}} = (1-\beta)\cdot(1-\tau_c)
\label{eq:rholoop}
\end{equation}
The combined closed-loop gain, accounting for the model's open-loop
amplification $g$, is:
\begin{equation}
\rho_{\text{closed}} = g \cdot \rho_{\text{loop}}
\end{equation}
Bounded drift requires $\rho_{\text{closed}} < 1$.  With $\beta=0.40$ and
$\tau_c=0.05$, $\rho_{\text{loop}} = 0.57$, providing $\approx43\%$ gain
reduction per chunk before accounting for $g$.

\paragraph{Why the two factors multiply.}  $(1-\tau_c)$ and $(1-\beta)$ act on
disjoint tensor regions at different points in the cycle --- $(1-\tau_c)$ on
the \emph{overlap} region when it is \emph{read} as context
(Section~\ref{sec:renoise}), $(1-\beta)$ on the \emph{new-frame} region when it
is \emph{written} as output (Section~\ref{sec:adain}).  Each attenuator sits
on the one path the error must traverse per loop iteration: a perturbation in
chunk $N$'s output is written with gain $(1-\beta)$ (AdaIN is applied uniformly
across $L_N^{\text{new}}$, including the tail frames that become the next
SLB), then read back with gain $(1-\tau_c)$ once it becomes chunk $N{+}1$'s
overlap context.  For a linear, time-invariant loop this is the standard
construction of a loop gain from attenuators in series.  What
Eq.~\eqref{eq:rholoop} does not capture is any cross-term where the correction
strength itself depends on chunk content (e.g.\ the AdaIN amplification cap
engages non-uniformly); that content-dependence is exactly why $g$ cannot be
folded into a single constant either (below).

\paragraph{Hyperparameter choices.}
$\beta=0.40$ is the largest value that preserves scene-change responsiveness in
practice: at $\beta\geq0.5$ slow colour shifts accumulate due to excessive EMA
anchoring.  $\tau_c=0.05$ is a conservative re-noising level that maintains
high boundary similarity ($>0.97$) while partially correcting SLB distribution
shift; larger values ($\tau_c\geq0.20$) cause visible motion discontinuities on
22B streaming.  Together they yield $\rho_{\text{loop}}=0.57$.
$\lambda=0.10$ (EMA rate) was chosen to respond to gradual scene evolution
without over-weighting individual chunks; $\boldsymbol\gamma=(0.0,\,0.05,\,0.2,\,0.3)$
are conservative exponents that attenuate only bands with clear excess energy.

\paragraph{Per-chunk gain measurement.}
$g$ is measured empirically by running a second denoising pass with a slightly
perturbed overlap latent $\tilde{L}_{\text{overlap}} + \delta$
($\|\delta\|_F = \epsilon\,\|\tilde{L}_{\text{overlap}}\|_F$, $\epsilon=0.01$)
and comparing the output perturbation to the input:
\begin{equation}
g_{\rm SLB} = \frac{\|\Delta_{\rm output}^{\rm SLB}\|_F}{\|\delta\|_F}
\label{eq:gslb}
\end{equation}
where $\Delta_{\rm output}^{\rm SLB}$ is restricted to the last
$\text{overlap\_lf}$ new frames---the portion that becomes the next chunk's
SLB.  This measures sensitivity of the boundary output to boundary input
perturbations, which is the relevant quantity for SLB-mediated drift
propagation.  It does not capture drift arising from other channels (e.g.,
prompt conditioning or attention to non-overlap frames), so $g_{\rm SLB}$
is a lower bound on the true open-loop gain.  The measurement is exposed on
the \textbf{CLSSConfig} node (\texttt{measure\_g=on}, $\epsilon$ via
\texttt{measure\_g\_epsilon}); it doubles Stage-1 compute for the
continuation chunks it runs on.  Measured values at both the 16-lf and the
production 52-lf geometry, and their consequences for the stability model, are
reported in Section~\ref{sec:discussion}.

%--------------------------------------------------------------------------
\section{Algorithm}
\label{sec:algorithm}
%--------------------------------------------------------------------------

Algorithm~\ref{alg:clss} summarises the per-chunk CLSS procedure.

\begin{algorithm}[H]
\caption{CLSS per-chunk generation}
\label{alg:clss}
\begin{algorithmic}[1]
\Require Config: $\tau_c, \beta, \lambda, \boldsymbol\gamma$; state: SLB, EMA, anchor bank
\State $L_{\text{ov}} \leftarrow$ SLB.read()  \Comment{may be empty for chunk 0}
\State $\tilde{L}_{\text{ov}} \leftarrow (1-\tau_c)\,L_{\text{ov}} + \tau_c\,\varepsilon$ \Comment{Eq.~\eqref{eq:renoise}; Sect.~\ref{sec:renoise}}
\State $A \leftarrow$ anchor\_bank.retrieve(top-$m$ by cosine sim) \Comment{Sect.~\ref{sec:anchor}}
\State $L_N \leftarrow$ Denoise($\tilde{L}_{\text{ov}},\, A,\, \text{prompt}_N$)
\State $L_N^{\rm new} \leftarrow L_N[\text{overlap\_lf}:]$ \Comment{drop overlap region}
\State $L_N^{\rm new} \leftarrow$ AdaIN\_blend($L_N^{\rm new}$, EMA, $\beta$) \Comment{Sect.~\ref{sec:adain}}
\State $L_N^{\rm new} \leftarrow$ band\_soft\_shrink($L_N^{\rm new}$, band\_EMA, $\boldsymbol\gamma$) \Comment{Sect.~\ref{sec:shrink}}
\State $L_N^{\rm new} \leftarrow$ detail\_anchor($L_N^{\rm new}$, $E^{\rm ref}_{\rm lo/hi}$) \Comment{Eq.~\eqref{eq:detail}; Sect.~\ref{sec:detail}}
\State EMA.update($L_N^{\rm new}$, $\lambda$); band\_EMA.update($L_N^{\rm new}$, $\lambda$)
\State anchor\_bank.update(last frame of $L_N^{\rm new}$) \Comment{Sect.~\ref{sec:anchor}}
\State SLB.push($L_N^{\rm new}[-\text{overlap\_lf}:]$)
\State \Return $L_N^{\rm new}$
\end{algorithmic}
\end{algorithm}

%--------------------------------------------------------------------------
\section{Implementation}
\label{sec:impl}
%--------------------------------------------------------------------------

\subsection{Memory Budget for \texorpdfstring{16\,GB}{16 GB} VRAM}
\label{sec:memory}

LTX-2.3 22B in BF16 occupies $\approx44$\,GB.  Three complementary strategies
fit it within 16\,GB VRAM:

\begin{enumerate}[leftmargin=*,noitemsep]
\item \textbf{GGUF 4-bit quantisation} (Q4\_K\_S): reduces the on-disk size
  to $\approx11$\,GB.  Dequantization to BF16 at load time requires
  $\approx44$\,GB of pinned CPU RAM.
\item \textbf{CPU block-streaming}: all transformer blocks reside in pinned
  CPU RAM; only a small number of blocks are resident on GPU at a time, with
  the next block pre-fetched asynchronously via a dedicated CUDA stream.
\item \textbf{Temporal chunking (CLSS)}: latent memory is $O(\text{overlap\_lf})$
  instead of $O(\text{total\_frames})$.  Processed chunk latents are immediately
  offloaded to CPU after each chunk.
\end{enumerate}

Combined peak VRAM budget at the production resolution:
\begin{center}
\begin{tabular}{lr}
\toprule
Component & Peak VRAM \\
\midrule
Model blocks (BF16, streaming) & $\approx5$\,GB \\
Activations + latents (per chunk) & $\approx4$\,GB \\
VAE / text encoder (loaded/freed) & $\approx2$\,GB \\
\midrule
\textbf{Total peak} & $\approx11$\,GB \\
\bottomrule
\end{tabular}
\end{center}

\subsection{Two-Stage Pipeline}
\label{sec:twostage}

\paragraph{Stage 1 (Open-resolution streaming).}
CLSS streaming generates the full video at a coarse latent resolution
($H_L \times W_L$ latent grid, $11\times20$ in the production runs).  Each
chunk runs the full 20-step LTX scheduler.  CLSS corrections are applied
between chunks as per Algorithm~\ref{alg:clss}.

\paragraph{Stage 2 (Distilled refinement).}
Stage 1's coarse resolution and 20-step schedule are sufficient for temporal
structure and scene content but produce visible quantisation artefacts and
reduced texture detail after GGUF 4-bit dequantisation.  Stage 2 addresses
this by upsampling and refining the output.

The Stage 1 output is spatially upsampled $2\times$ by
\texttt{LTXVLatentUpsampler}, yielding a $22\times40$ latent grid
($1408\times1024$ pixels).  Stage 2 then refines the full video in temporal
chunks using a distilled schedule with initial noise level
$\sigma_0 \approx 0.909$ and a dense manual sigma sequence
($0.9094,\,0.725,\,0.55,\,0.4219,\,0.25,\,0$; five steps).  Each chunk's new
frames are seeded from the upscaled Stage 1 latent with $\sigma_0$-scaled
noise:
\begin{equation}
x_{\text{start}} = \sigma_0\,\varepsilon + (1-\sigma_0)\,L_N^{\text{S1}}
\end{equation}
anchoring Stage 2 to the coarse temporal structure produced by Stage 1 while
allowing the refinement to add local high-frequency detail.  The chunk
geometry is auto-selected from a VRAM token budget (production runs:
$\text{frames\_per\_chunk}=42$, $\text{overlap}=14$).

Key differences from Stage 1:
\begin{itemize}[noitemsep]
\item AdaIN and spectral-shrinkage corrections are \emph{disabled}: Stage 2
  refinement is structurally anchored to the S1 latent, so distributional
  drift correction is unnecessary.
\item The detail anchor (Section~\ref{sec:detail}) \emph{does} run, per
  window, with the reference recomputed from that window's upscaled Stage-1
  slice rather than carried from a first-window EMA: the distilled refinement
  systematically adds spatial high-frequency energy relative to its Stage-1
  seed, and per-window re-referencing means the correction cannot accumulate.
\item A global pre-generated noise field is sliced per chunk
  (\texttt{\_SlicedNoise}) so the per-chunk noise samples tile seamlessly at
  the chunk boundaries.
\item Audio is refined under a join-sigma rule
  (Section~\ref{sec:s2audiorefine}).
\end{itemize}

\paragraph{Stage-2 audio refinement with a join sigma.}
\label{sec:s2audiorefine}
Freezing the audio through Stage~2 leaves a structural asymmetry: the video
receives a dense low-$\sigma$ descent (re-noised to $\sigma_0\approx0.909$
and re-denoised over the distilled steps) while the audio would receive none.
Simply unfreezing does not work --- and the reason is informative.  Full
refinement from $\sigma_0\approx0.909$ was measured to destroy the audio
(\texttt{aud\_refine\_sim} $0.07$--$0.16$: only $7$--$16\%$ of the Stage-1
audio survives, each window re-inventing its own content), because a melody,
timbre, or voice is not in the text prompt --- re-inventing audio from $91\%$
noise means losing it, whereas the video survives the same treatment because
its composition \emph{is} prompt-specified and the detail anchor re-anchors
it.  The fix is a \emph{join sigma} (\texttt{audio\_join\_sigma}): the audio
stays frozen (mask 0, audio timestep 0) through every step with
$\sigma > \sigma_{\text{join}}$ and joins the descent at the first step with
$\sigma \le \sigma_{\text{join}}$, with an explicit noise-up to that
$\sigma$, keeping a fraction $(1-\sigma_{\text{join}})$ of the Stage-1 audio
as starting signal.  The descent is $\sigma$-matched at every step --- both
modalities sit at the current step's sigma, so the failure mode of a
per-modality schedule split (Section~\ref{sec:sigmasplit}) does not arise ---
and it costs no extra compute, since Stage~2 already pays for these passes
and the audio is already in the (masked) latent.  The production
configuration uses $\sigma_{\text{join}}=0.55$ on the dense schedule ($45\%$
of the Stage-1 audio kept, three matched steps); measured results are
reported in Section~\ref{sec:exp_join}.  A chunk-1 onset limiter (smooth
$3.5\sigma\to4\sigma$ tanh) is applied to the refined audio because the
chunk-1 re-roll otherwise re-creates a violent onset transient (measured
$-6.08$, $\approx11\sigma$ at the run's std).  The known remaining weakness
is content seams at the Stage-2 window boundaries (\texttt{aud\_bnd}
$0.30$--$0.61$ across the production runs).

\subsection{Audio Streaming}
\label{sec:audio}

Audio and video share the same temporal chunk boundaries.  Audio length
accounting is anchored to \emph{pixel time}: the causal video VAE maps a
chunk of $\text{new\_lf}$ latent frames to $(\text{new\_lf}-1)\cdot 8+1$
pixel frames for the first chunk of a sequence and $\text{new\_lf}\cdot 8$
pixel frames for every continuation chunk.  The audio rate is taken from the
chunk template, $r = \text{new\_af} / \big((\text{new\_lf}-1)\cdot 8+1\big)$
audio frames per pixel frame, and the SLB length follows the same rate:
\begin{equation}
\text{overlap\_af} = \text{round}\!\left(\text{overlap\_lf}\cdot 8 \cdot r\right)
\end{equation}
In the production geometry ($52$\,lf new, $8$\,lf overlap) this gives
$\text{new\_af}=427$ (chunk 1) / $434$ (continuations), $\text{overlap\_af}=67$.
The effective video overlap is additionally clamped to $\text{new\_lf}$: the
SLB can never carry more frames than a chunk produces, and unclamped values
were measured to silently discard motion at every seam.

\paragraph{Continuity via reference audio and the SLB.}
Audio continuity is carried by two conditioning paths.  The overlap-time audio
of each continuation chunk carries an audio analogue of the video SLB: the
overlap audio frames are seeded from the previous chunk's tail and re-noised
at $\tau_c$, so each chunk opens with near-verbatim continuity.  Alongside
the SLB, a reference-audio window (\texttt{ref\_audio} tokens at negative
temporal RoPE) taken from the previous chunk's pre-overlap tail provides
content conditioning; its length is \texttt{ref\_audio\_seconds} (production
value $8.0$\,s $=200$ frames).  The reference default is the audio-overlap
length ($67$ frames $\approx2.7$\,s).  The audio context (SLB plus reference)
is envelope-flattened before conditioning: each context frame is rescaled
toward the window-mean RMS (gain capped at $[0.6,\,1.25]$, asymmetric so quiet
frames are never over-amplified), removing the loudness arc while preserving
content and timbre.

The audio SLB is re-noised on its own schedule.  At the production geometry
the multiplier is $1.0$ (the base $\tau_c$ schedule with an audio ceiling of
$0.35$); at shorter chunks ($16$\,lf) a hotter $3\times$ schedule was
perceptually validated against a per-chunk loudness-gesture repetition: the near-verbatim SLB tail circulates the previous chunk's
loudness envelope, and hotter re-noising breaks the loop while the
envelope-flattened context removes the arc from the conditioning.

\paragraph{Upward-only RMS anchoring.}
Un-anchored streaming audio is a divergent loop: runs were measured with RMS
growing $+67\%$ and low-mid spectral bands tripling over ten chunks.  A
per-chunk anchor corrects the new audio's global RMS toward an EMA target
whose drift from the chunk-1 reference (computed onset-excluded) is capped at
$\pm5\%$, before the chunk's tail is fed forward as the next reference
window.  The anchor direction is \emph{upward-only}, matching the reference
streaming pipeline: a chunk quieter than the capped EMA target is boosted
toward it, while a louder chunk is kept raw.  A hard bidirectional version
was measured cutting brighter chunks back toward the chunk-1 level --- a
dullness-lock.  The final-assembly spectral normalisation is likewise
upward-only (per-bin gains floored at $1.0$, RMS pin boosts only), preserving
the anti-drone direction (a high-frequency decay being boosted back) while a
genuine low-frequency rise passes through and is gated on the
\texttt{aud\_hf} trend.  One caveat is explicit: the loudness runaway the
upward-only law cannot suppress (RMS $\times4$ over a 15-chunk run measured
earlier) remains unhandled, so long runs gate on the \texttt{aud\_rms} trend.

\paragraph{Boundary smoothing.}
At each audio chunk boundary, a linear blend over $\pm2$ audio latent frames
is applied in latent space before VAE decoding.  This eliminates click
artifacts that arise from the abrupt splice between consecutive chunks,
which would otherwise be visible as discontinuities in the spectrogram.

\subsection{Per-Modality Sigma Schedules}
\label{sec:sigmasplit}

The noise schedule is not a free parameter in this pipeline: ComfyUI's
\texttt{LTXVScheduler} derives its sigma shift from the token count of the
latent template connected to it,
\begin{equation}
\sigma_{\text{shift}} = \text{tokens}\cdot\frac{1.1}{3072} + 0.583,
\qquad \text{tokens} = \text{template\_lf}\times H\times W ,
\label{eq:sigmashift}
\end{equation}
so the schedule changes whenever the chunk length does.  At the
$11\times20$ production spatial template a 16-lf chunk template yields
shift $1.844$ and a 52-lf template yields $4.680$; the formula reproduces
the logged sigma schedules to a maximum absolute error of $5\times10^{-5}$.
This matters because the two modalities respond very differently to the
schedule at the production chunk length (52\,lf $\approx17.3$\,s windows),
and the conflict was measured in both directions:

\begin{itemize}[leftmargin=*,noitemsep]
\item \textbf{Shared 4.68 schedule} (13 of 20 steps at $\sigma\ge0.9$):
  the video is well-formed, but the audio is under-developed --- the raw
  Stage-1 audio latent lands at std $\approx0.42$ against $0.98$--$1.26$
  for real music through the same VAE (Section~\ref{sec:exp_scale}), and
  the decoded output measures $-34$\,dB RMS, L/R correlation
  $0.9988$--$0.9994$ (stereo collapsed to mono), ro99 bandwidth
  $4.4$--$4.7$\,kHz (real music: $6$--$12$\,kHz).
\item \textbf{Shared flat 1.844 schedule at 52-lf chunks}: the audio is
  restored (raw std $0.49$, validated perceptually) but the video is destroyed
  --- it collapses to a near-static copy of the guide image (video latent
  std $1.99$ against a normal $\approx1.01$; origin similarity to the guide
  frame $\to0.999$ by frame~3, i.e.\ every frame \emph{is} the guide image).
\end{itemize}

At 52-lf chunks the video wants the top-heavy 4.68 schedule and the audio
wants the 1.844 schedule, and no shared schedule serves both.  The
resolution is a \emph{per-modality sigma schedule inside the single joint
pass}: the video follows the original LTXVScheduler sigmas byte-identical
(its ancestral noise stream is untouched --- the audio draws from a separate
RNG stream), while the audio follows the schedule generated by
Eq.~\eqref{eq:sigmashift} with $\sigma_{\text{shift}}$ scaled by
$\text{mult} = 1.844/4.680 \approx 0.394$, keeping $\sigma_0 = 1.0$: the
audio still starts from full noise and only the \emph{shape} of its
schedule is remapped.  Two mechanisms make the split exact rather than
approximate: a \texttt{process\_timestep} patch reports the audio's true
per-token sigma ($\text{mask}\times a_i$) to the model, and because
\texttt{calculate\_denoised} converts the velocity prediction with the
\emph{global} sigma, the sampler re-derives the audio $x_0$ with $a_i$
instead --- mask-aware, so only fully-denoised audio tokens are converted
and mask-pinned SLB tokens are handled by the mask itself.  The multiplier
is auto-derived (\texttt{audio\_shift\_mult}$=0$): the node divides the
perceptually validated audio target shift $1.844$ by the video shift it recovers
from the connected sigmas.  Live validation is reported in
Section~\ref{sec:exp_split}.

\subsection{ComfyUI Node Integration}
\label{sec:nodes}

Six custom nodes implement the full pipeline inside ComfyUI:

\begin{enumerate}[leftmargin=*,noitemsep]
\item \textbf{CLSSConfig}: exposes the CLSS hyperparameters worth exposing
  ($\tau_c, \beta$, overlap); validated internals are fixed.  Diagnostic
  levers: \texttt{shrinkage=off} zeroes $\boldsymbol\gamma$ (the
  spectral-shrinkage ablation arm), and \texttt{measure\_g} (with
  \texttt{measure\_g\_epsilon}) enables the $g_{\rm SLB}$ measurement of
  Eq.~\eqref{eq:gslb}.
\item \textbf{CLSSScenePrompts}: multi-scene text input via \texttt{---}
  delimiters; Gemma-enhanced per scene.  Output is a flat CONDITIONING list
  with one entry per scene, proportionally assigned to chunks.
\item \textbf{CLSSStreamingSampler}: Stage 1 streaming sampler.  Accepts any
  ComfyUI GUIDER/SAMPLER/SIGMAS/NOISE/LATENT combination.  Emits an
  unconditional settings dump, per-chunk telemetry, end-of-run trend
  summaries with per-chunk localisation, per-frame origin/layout drift event
  tables, and a phase-lock repetition detector.  Experiment levers:
  \texttt{open\_loop=on} disables every closed-loop correction to give the
  naive sliding-window baseline arm; \texttt{object\_identity=on} adds the
  spatially-resolved seam identity metric; \texttt{auto\_seam\_pin=on}
  enables the event-triggered seam pin of Section~\ref{sec:events}.
\item \textbf{CLSSStage2}: Stage 2 refinement sampler with consistent per-chunk
  noise slicing, per-window detail anchoring, and the joined audio refinement
  of Section~\ref{sec:s2audiorefine}.
\item \textbf{CLSSAVGuiderV2}: the multi-modal guider (Section~\ref{sec:guider}).
\item \textbf{CLSSAVGuider}: the legacy split-CFG patch over an existing
  guider, retained for compatibility and inert on current ComfyUI.
\end{enumerate}

All stochastic inputs are seeded: Stage-1 video \emph{and} audio noise are
pre-generated as full-sequence fields sliced per chunk (distinct yet coherent
noise per chunk, bit-reproducible runs), and every fallback noise path uses a
generator derived deterministically from the seed and chunk position.

\subsection{Multi-Modal Guidance and the Packed-Latent Pitfall}
\label{sec:guider}

The reference pipeline guides each modality independently:
\begin{equation}
\text{pred} = c + (s_{\text{cfg}}-1)(c-u)
             + s_{\text{stg}}\,(c-p)
             + (s_{\text{mod}}-1)(c-m)
\label{eq:guider}
\end{equation}
per modality, followed by a rescale
$\text{pred} \mathrel{*}= r\,\sigma_c/\sigma_{\text{pred}} + (1-r)$,
where $u$ is the text-unconditional prediction, $m$ the modality-isolated
prediction (both audio$\leftrightarrow$video cross-attentions skipped), and
$p$ the STG-perturbed prediction.  The STG perturbation skips
\emph{self}-attention in designated blocks (block 28 for LTX-2.3) by
replacing the attention output with the raw value projection
($\text{out}=W_V x$, gating and output projection unchanged) --- a
value-passthrough, not a zero-out.  The working configuration runs
per-modality CFG $4.0$ video / $7.0$ audio, STG $1.0$ video / $0$ audio,
$s_{\text{mod}}=1$ (the modality term of Eq.~\eqref{eq:guider} vanishes and
its pass is skipped --- three passes per step), rescale $0.7$.  Audio STG is
off: a joint run with audio STG on deflated the audio, and every
perceptually validated run used audio STG $0$.

\paragraph{The packed-latent representation.}  ComfyUI's
\texttt{CFGGuider.sample()} packs nested audio-video latents into a
\emph{single flat} $[B,1,N_v{+}N_a]$ tensor before sampling and unpacks only
after.  A guider that tries to detect the audio-video case with
\texttt{isinstance(x, NestedTensor)} inside \texttt{predict\_noise} therefore
never sees a nested tensor and falls back to plain shared CFG --- and it does
so silently.  The guider must instead capture the per-modality shapes in
\texttt{sample()} \emph{before} packing, detect the packed $[B,1,N]$
representation in \texttt{predict\_noise}, unpack each pass output, apply
Eq.~\eqref{eq:guider} per modality, and repack.  Two diagnostics guard
against a silent regression: the guider announces the active branch and its
pass count once per chunk, and it logs the STG difference norms so a
zero-valued (inert) perturbation is visible immediately.  A useful invariant
follows from the pass count: because each guidance term costs a transformer
pass, a guidance configuration that does not change per-step wall-clock time
is not actually running.

%--------------------------------------------------------------------------
\section{Experiments}
\label{sec:experiments}
%--------------------------------------------------------------------------

All runs use the production configuration of Table~\ref{tab:config}: a
single prompted scene, streamed as four chunks of $52$ latent frames
($\approx17.3$\,s each) with $8$-frame overlap, assembling to $208$ latent
frames ($69.3$\,s at 24\,fps), refined by a five-chunk Stage~2 pass at
$2\times$ resolution.  The telemetry cited below is the per-chunk output of
the pipeline itself (settings dumps, per-chunk metrics, trend summaries,
event tables); the final decoded outputs of the production I2V and T2V runs
were verified perceptually by the authors.
\label{sec:telemetry}

\subsection{Shared Configuration}

\begin{table}[H]
\centering
\caption{Production configuration shared by all runs.  Guidance values are
those actually active during sampling (verified by the guider's per-chunk
\texttt{passes/step} and difference-norm telemetry, Section~\ref{sec:guider}).}
\label{tab:config}
\begin{tabular}{ll}
\toprule
Parameter & Value \\
\midrule
Model & LTX-2.3 22B (Q4\_K\_S GGUF, BF16 CPU dequant) \\
Geometry & 4 chunks $\times$ 52 lf new $+$ 8 lf overlap $=208$ lf ($69.3$\,s) \\
Resolution & $704\times1280$ (Stage 1, $11\times20$ latent); $2\times$ (Stage 2) \\
$\tau_c$ / $\beta$ / $\lambda$ / $\Delta_\sigma$ & 0.05 / 0.40 / 0.10 / 0.05 \\
$\boldsymbol\gamma$ & $(0.0,\,0.05,\,0.2,\,0.3)$ \\
Anchor bank & threshold 0.78, persistence 2, max 8, top-2; $F_a$ auto (=2) \\
Detail anchor & on (Stage 1 scene-first ref; Stage 2 per-window ref) \\
Video / audio guidance & CFG 4.0 / 7.0; STG 1.0 / 0; modality 1.0; rescale 0.70 \\
STG block & 28 (three passes per step) \\
Audio sigma split & auto: $\text{mult}=1.844/4.680=0.394$ (shape-only, full-noise start) \\
Audio continuity & SLB 67\,af re-noised at $\tau_c$ (ceiling 0.35); ref 200\,af (8.0\,s), negative RoPE \\
Audio anchor & RMS-only, upward-only, $\pm5\%$ EMA cap \\
S1 noise temporal-corr mix & 0.30 (video only) \\
Stage 1 steps & 20 (LTXVScheduler, connected template) \\
Stage 2 & dense sigmas $0.9094/0.725/0.55/0.4219/0.25/0$, audio join $\sigma=0.55$, \texttt{audio\_refine=on} \\
Stage 2 geometry & auto: $\text{fpc}=42$, overlap $14$ ($T=208$, $22\times40$) \\
Seeds & Stage 1 $1125899906842624$; Stage 2 $809542877813341$ \\
Diagnostics & \texttt{object\_identity=on}; \texttt{open\_loop}/\texttt{auto\_seam\_pin} per experiment \\
VRAM & 15.61\,GB GPU; peak usage within budget \\
\midrule
T2V first frame & text only \\
I2V first frame & guide image, latent $[1,128,1,11,20]$ \\
\bottomrule
\end{tabular}
\end{table}

%--------------------------------------------------------------------------
\subsection{Closed-Loop versus Open-Loop Streaming}
\label{sec:exp_ab}
%--------------------------------------------------------------------------

The exposure-bias account (Section~\ref{sec:background}) predicts that
removing the corrections at the production geometry re-opens drift.  Two
same-seed I2V runs test this: \texttt{open\_loop=on} (every closed-loop
correction disabled: re-noising, AdaIN, shrinkage, detail and std anchors,
audio RMS anchor) over five chunks ($86.7$\,s), and the closed-loop arm over
four chunks ($69.3$\,s), both with \texttt{measure\_g=on}.

Without corrections, drift returns exactly as predicted: video std
$1.015\to1.133$ ($+0.118$ over five chunks), audio RMS $0.376\to0.504$
($+34\%$), audio high-frequency energy drifting to $1.60\times$ the chunk-1
reference (chunk-5 per-bin ratios up to $2.09$), and the loudness-envelope
correlation between consecutive chunks rising $-0.05\to+0.59$.  The
closed-loop arm holds the lines the open arm loses: video std $+0.001$,
object identity $0.980$--$0.995$ (per-cell minimum $0.79$ at the chunk-2
seam), audio high-frequency drift inside a $0.22$ range.  Audio RMS rises
$+0.087$ --- partly by design: the upward-only anchor lets louder chunks keep
their gain.  The measured open-loop gain (Eq.~\eqref{eq:gslb}) reads
$g_{\rm SLB}=24.8/28.8/42.4$ on the closed arm (chunks 2--4) and
$33.8/35.9/38.4/21.4$ on the open arm (chunks 2--5) --- the same $15$--$87$
range measured at the shorter 16-lf geometry, so the transformer's boundary
sensitivity does not change with chunk length.  The arms share seed and
configuration but differ in length (4 vs.\ 5 chunks), and same-seed runs at
different total lengths do not share a prefix (the pre-generated noise layout
depends on total length): close to controlled, not perfectly.

%--------------------------------------------------------------------------
\subsection{Image-to-Video Generation (\texorpdfstring{69.3\,s}{69.3 s})}
\label{sec:exp_i2v}
%--------------------------------------------------------------------------

The I2V run (guide image, four 52-lf chunks) is the reference point of the
paper; its decoded output was verified perceptually.

\paragraph{Stage 1 video.}  Video latent std holds $1.011\to1.006$
($\Delta-0.004$).  Boundary similarity across the three seams reads
$0.949/0.977/0.986$; intra-chunk similarity grows $0.692\to0.910$; identity
to the anchor bank holds $0.949\to0.986$; the decoded seam-grid object
identity (8$\times$8 cell cosine) holds mean $0.988/0.996/0.998$ with
per-cell minima $0.875/0.963/0.988$.  The detail anchor holds the
high-frequency energy share flat end-to-end ($0.527\to0.527$).  Origin
similarity to the scene-first frame walks $0.378\to0.331$ --- the scene
moving away from its opening, not a drift failure.

\paragraph{Stage 1 audio.}  Raw audio latent std per chunk reads
$0.440/0.400/0.457/0.495$ --- the split schedule
(Section~\ref{sec:sigmasplit}) restoring the audio scale.  The RMS anchor
(upward-only) holds loudness to the chunk-1 reference:
$0.435\to0.504$ over the run, with two chunks kept raw (louder, by design)
and one boosted ($g=1.09$).  The audio SLB is honoured exactly on every
continuation chunk ($1.0$) despite the re-noising; the loudness-envelope
correlation across seams reads $0.032\to0.477\to0.808$ --- the climb at this
chunk length is a property of the continuity context's shared envelope, and
the output was verified free of audible repetition.  Audio boundary
similarity reads $0.624/0.295/0.357$ (the chunk-3 seam is the run's weakest
audio point); audio high-frequency share rises $0.775\to1.067$ inside the
stability tolerance.

\paragraph{Stage 2.}  Five windows (auto geometry, Table~\ref{tab:config}).
Video fidelity to the Stage-1 latent holds $0.981\to0.893$ (first frame) and
$0.958\to0.882$ (last frame) --- the known late-run softening.  The joined
audio refinement (Section~\ref{sec:s2audiorefine}) holds
\texttt{aud\_refine\_sim} $0.66/0.67/0.69/0.73/0.74$ across the windows and
leaves the final audio at std $0.517$ against Stage-1's $0.459$: brighter,
not re-invented.  The audio seams between Stage-2 windows read
$0.43/0.39/0.44/0.61$ --- the joined descent's known residual.

%--------------------------------------------------------------------------
\subsection{Text-to-Video Generation (\texorpdfstring{69.3\,s}{69.3 s})}
\label{sec:exp_t2v}
%--------------------------------------------------------------------------

The T2V run (text only, same seed and configuration; the seam-pin safeguard
was enabled and did not trigger) was likewise verified perceptually.

\paragraph{Stage 1 video.}  Video latent std holds $1.036\to1.038$
($\Delta+0.002$).  Boundary similarity reads $0.929/0.829/0.978$;
intra-chunk similarity grows $0.784\to0.883$; anchor identity holds
$0.929\to0.978$; seam-grid object identity holds $0.997\to0.994$.  The
detail anchor holds the high-frequency share flat ($0.616\to0.616$).  Origin
similarity walks $0.612\to0.476$ (flagged drift) --- without a guide image
the text-only scene moves further from its opening than the I2V run does,
while the anchored identity metrics stay put.

\paragraph{Stage 1 audio.}  RMS $0.446\to0.432$ against the chunk-1
reference; the audio SLB honoured exactly on every continuation chunk;
loudness-envelope correlation across seams reads $-0.096\to0.555$, below the
repetition flag at every seam.  Audio boundary similarity reads
$0.361/0.246/0.198$; high-frequency share falls $1.198\to1.074$.

\paragraph{Stage 2.}  Video fidelity to Stage-1 holds $0.921\to0.826$
(first) and $0.920\to0.828$ (last).  Joined audio refinement holds
\texttt{aud\_refine\_sim} $0.58/0.65/0.64/0.67/0.65$; audio seams between
Stage-2 windows read $0.53/0.52/0.39/0.45$.

%--------------------------------------------------------------------------
\subsection{The Per-Modality Schedule Split}
\label{sec:exp_split}
%--------------------------------------------------------------------------

The split of Section~\ref{sec:sigmasplit} was validated before it became the
production default (I2V, $2\times52$\,lf): the video runs the byte-identical
4.68-shift sigmas and the audio its own $0.39\times$-shifted schedule
($\sigma_{\text{shift}}\approx1.83$, full-noise start),
\texttt{audio\_stg}$=0$.  The audio recovers the scale the shared 4.68
schedule suppressed --- raw Stage-1 audio latent std $0.497$ against $0.42$
on the shared schedule --- while the video is untouched: video std
$1.013\to1.008$, boundary similarity $0.967$, identity $0.967$, all in the
normal range on the unchanged video sigmas.  Stage-2 audio refinement
(join sigma) holds \texttt{aud\_refine\_sim} $0.70$--$0.72$, and the final
audio std is $0.540$.  The output was validated perceptually.  In the
production runs the multiplier is auto-derived (\texttt{audio\_shift\_mult}$=0$):
$1.844/4.680=0.394$, applied shape-only with a full-noise start.  An earlier
joint run of the same split deflated the audio (std $0.30$), but that run
also carried \texttt{audio\_stg}$=1$ on the guider --- every perceptually
validated run used \texttt{audio\_stg}$=0$ --- so the deflation is
attributed to the guidance confound, not the split.

%--------------------------------------------------------------------------
\subsection{Joined Stage-2 Audio Refinement}
\label{sec:exp_join}
%--------------------------------------------------------------------------

The join-sigma design of Section~\ref{sec:s2audiorefine} is the outcome of a
measured escalation.  Full refinement from $\sigma_0=0.909$ destroyed the
audio (\texttt{aud\_refine\_sim} $0.07$--$0.16$): a melody, timbre, or voice
is not in the text prompt, so re-inventing it from $91\%$ noise loses it.
Intermediate joins trace the boundary between the regimes: a join at
$0.725$ (stock three-step sigmas) held $0.45$--$0.49$; a join at $0.422$
kept $58\%$ of the Stage-1 audio and read high fidelity but without the
brightness gain (rejected); the production dense schedule with
$\sigma_{\text{join}}=0.55$ keeps $45\%$ of the Stage-1 audio as starting
signal over three matched steps and holds \texttt{aud\_refine\_sim}
$0.66$--$0.74$ in the I2V run (Section~\ref{sec:exp_i2v}) and $0.58$--$0.67$
in the T2V run (Section~\ref{sec:exp_t2v}), brightening the final audio (std $0.459\to0.517$,
I2V).  The chunk-1 onset limiter is load-bearing: without it the chunk-1
re-roll re-creates a measured $-6.08$ ($\approx11\sigma$) onset transient.
The rule the escalation settled: audio is either generated in one pass, or
re-denoised only from a noise level low enough that its content survives.

%--------------------------------------------------------------------------
\subsection{Audio Latent Scale and the Measurement Probe}
\label{sec:exp_scale}
%--------------------------------------------------------------------------

The audio VAE's latent normaliser is calibrated so that \emph{real} audio
lands at unit variance: encoding four real music tracks gives an overall
latent std of $0.98$--$1.26$ with the per-frequency-bin std flat at $1.00$
($0.94$--$1.07$ across all 16 bins).  Generated audio does not get there, and
the shortfall predicts the perceptual failure: the half-scale latents decode
to exactly what is heard --- quiet, low-passed near $3$--$4$\,kHz against the
$8$\,kHz mel ceiling, and, at the longest windows, stereo collapsed to
identical channels.

\begin{table}[H]
\centering
\caption{Latent statistics of real and generated audio through the same VAE
(probe: \texttt{simulations/audio\_latent\_probe.py}).  ro99 is the 99th
percentile frequency; L/R is the inter-channel correlation of the latent.
Real music sits at unit variance; generated audio sits below it, and the
shortfall tracks the audible deficits.}
\label{tab:scale}
\begin{tabular}{lccc}
\toprule
Source & latent std & ro99 (kHz) & L/R corr \\
\midrule
Real music (4 tracks) & $0.98$--$1.26$ & $8$--$12$ & $0.69$--$0.87$ \\
Generated, 26\,s per-chunk & $0.84$ & $6.1$ & $0.93$ \\
Generated, 20\,s per-chunk & $0.94$ & $4.0$ & $0.98$ \\
Generated, single 106\,s window & $0.63$ & $3.7$ & $0.999$ \\
Generated, single 104\,s window & $0.48$ & $3.1$ & $1.00000$ \\
\bottomrule
\end{tabular}
\end{table}

The trade the table exposes: the single-window soundtrack draft buys global
planning but the two longest-window runs are also the two most
mean-regressed, while the strongest runs on these measures came from
per-chunk windows --- a mid-length window is the untested middle ground.
Contrast is \emph{not} the deficit: the draft's 0.25\,s envelope range
(8.1\,dB) already exceeds real music's (4.3--6.3\,dB) and its log-mel
contrast (1.35) exceeds real music's (1.00); only level and treble are
short.  The VAE itself is innocent: real music round-trips through it with
bandwidth and stereo intact.

The probe is a small CPU script (audio VAE only, no transformer --- it runs
while the GPU is busy) reporting latent std, envelope dynamic range (dynR),
L/R correlation, and ro99 bandwidth.  Of these, dynR is the one
perceptually validated axis: in a controlled A/B a $19.7$\,s window at dynR
$9.8$\,dB was judged much better than a $104.3$\,s window at $3.3$\,dB (a
drone), while bandwidth did \emph{not} separate the two --- bandwidth and
level are not quality, and optimising them is not optimising what is heard.
A post-hoc statistical repair confirms this: shifting the audio latent's
per-(channel,bin) mean toward a real-music profile improved every metric
being watched (ro99 $3141\to4477$\,Hz, level $+4$\,dB) and the live run was
noise --- the recovered ``bandwidth'' was a broadband noise bed and the
loudness envelope collapsed.  Latent metrics localise failures; they never
prove a quality win --- the ear is the only ground truth.

%--------------------------------------------------------------------------
\subsection{Layout Stability and the Seam-Pin Safeguard}
\label{sec:events}
%--------------------------------------------------------------------------

Coarse scene layout is monitored per frame by a sign-free power-map metric:
the channel-mean spatial map of each frame is squared, hard low-passed
($3\times3$, stride 3), and compared against the scene-first frame by cosine
similarity.  Because the maps are non-negative, the metric cannot flip sign
on fine latent phase, and the hard low-pass retains only the spatial scale
that survives decoding.  In the verified production outputs no mid-run
layout events are observed: layout similarity fluctuates smoothly within its
scene-normal range, and the sudden single-frame re-layouts the metric
exists to catch (visible as jumping or morphing content) do not occur in the
runs reported here.

As a safeguard against such events, the sampler exposes an event-triggered
seam pin (\texttt{auto\_seam\_pin}): when a chunk's layout series crosses
the event thresholds (maximum per-frame drop $>0.25$ or minimum $<0.05$),
the \emph{next} chunk's overlap is placed frozen ($\tau_c=0$), so
generation continues from the last good frames exactly --- the fresh SLB,
$0.3$\,s old, rather than the anchor bank, whose entries are one to two
chunks old.  The trigger is evaluated from chunk 2 onward (chunk 1 is scene
establishment) and cannot protect the final chunk.  The mechanism localises
events and anchors the continuation to the most recent frames; because the
runs reported here contain no qualifying event, its behaviour against a
genuine re-layout is left to future measurement, together with the
underlying question of how to stop a joint audio-video model from
re-laying-out the opening frames of a new chunk --- the same design problem
that blocks anchor injection (Section~\ref{sec:anchor}): conditioning on a
temporally-displaced keyframe without corrupting audio or morphing video
has no known safe wiring.

%--------------------------------------------------------------------------
\subsection{Computational Overhead}
\label{sec:cost}
%--------------------------------------------------------------------------

The corrections themselves add negligible compute: AdaIN, shrinkage, the
detail anchor, and the boundary bookkeeping operate on chunk-sized tensors
between denoising runs.  The dominant cost drivers are structural.  Stage 1
runs three transformer passes per step (conditional, unconditional,
STG-perturbed); measured at the 16-lf geometry, chunk 1 costs
$\approx20.6$\,s/it and continuation chunks (carrying the 8-frame overlap)
$\approx26.8$--$33.7$\,s/it for a $\approx53$-second clip, totalling
$\approx116$ minutes including the $\approx21$-minute Stage 2.  The
per-modality audio split, the joined Stage-2 refinement, and the
\texttt{object\_identity} decoding add no extra denoising passes; only
\texttt{measure\_g} does, doubling Stage-1 compute on the chunks it
measures.  In Stage 2 the per-window cost rises $\approx50\%$ after the
first window against a $\approx32\%$ token-count rise: as the run proceeds
an increasing share of the model is pushed to CPU, so the slowdown is an
offload-pressure effect rather than a compute-scaling one.  The production
52-lf geometry trades fewer seams (four chunks instead of ten) for larger
windows; the sampler logs a startup warning whenever the connected
template's token count pushes the schedule outside the validated shift
range, so a silently changed schedule cannot go unnoticed.

%--------------------------------------------------------------------------
\section{Discussion}
\label{sec:discussion}
%--------------------------------------------------------------------------

\subsection{What the Closed Loop Bounds, and What It Does Not}
The open/closed ablation (Section~\ref{sec:exp_ab}) is the cleanest statement
of the loop's value: at the production geometry, every correction off
re-opens drift on every measured channel (video std $+0.118$, audio RMS
$+34\%$, audio HF $\to1.60\times$ reference, envelope correlation climbing),
and every correction on holds video std to $+0.001$ with object identity
$0.980$--$0.995$.  What the loop does \emph{not} bound, at either modality:
the loudness-envelope correlation between consecutive chunks still climbs at
52-lf chunks ($0.808$ by chunk 4 in the production I2V run, below the flag
in T2V), and the audio boundary metric floors mid-run in both production
runs.  The output-level verification matters here: the production I2V run
was verified free of audible repetition despite the envelope climb --- at
this chunk length the climb tracks the continuity context's shared envelope
rather than a replayed gesture.

\subsection{The Stability Model as a Local Design Guide}
\label{sec:stabdisc}
The measured boundary sensitivity $g_{\rm SLB}\approx15$--$87$ (16-lf
runs) and $21$--$42$ (52-lf runs) sits an order of magnitude above the
nominal stability boundary --- the naive product
$\rho_{\rm loop}\times g_{\rm SLB}\approx0.57\times26.7\approx15$ would
predict divergence, yet the closed-loop runs are observably stable (video
std $\Delta<0.005$, boundary similarity rising or flat).  The linear gain
model of Section~\ref{sec:stability} therefore does not describe the closed
loop asymptotically: the one-step sensitivity includes legitimate content
re-interpretation of a 1\% boundary perturbation, which the next denoise
pulls back toward the data manifold, and that contraction --- absent from
the model --- is what actually bounds drift.  The model remains useful as a
local design guide: which correction attenuates which channel, and by how
much (the product form of Eq.~\eqref{eq:rholoop} is validated by the
mechanism, not by the numbers).  A Lyapunov-style treatment or a Monte Carlo
sweep over perturbation seeds and chunk positions would be needed for a
rigorous account.

\subsection{Object Identity at the Seam}
The pooled identity metrics (mean cosine over the latent) are object-blind: a
vehicle or person can be replaced across a chunk boundary while the pooled
cosines stay above $0.95$.  The \texttt{object\_identity} metric closes the
observation gap by decoding the two seam frames and reporting a
spatially-resolved $8\times8$ grid-cell cosine, so a regional content swap
drops the minimum cell score even when the global mean stays high.  In both
production runs the metric corroborates the pooled telemetry: means
$0.988$--$0.998$ with per-cell minima $\ge0.875$ (I2V) and $\ge0.798$-era
pooled floors absent from the grid metric.  The mitigation for
object-critical content remains textual --- one prompt per chunk via
\textbf{CLSSScenePrompts}, re-describing the persistent objects explicitly.

\subsection{Scope and Generality}
\label{sec:scope}

\paragraph{Single configuration, single seed, single hardware.}
The production results share one set of hyperparameters, one seed, and one
hardware configuration (16\,GB VRAM); the T2V and I2V arms differ only in
the first-frame conditioning, by design.  Every single-run comparison is
therefore an $n=1$-per-arm observation, not a statistically powered result.
The ablation arms and the open-loop baseline are fully specified and exposed
as node levers (Section~\ref{sec:experiments}).

\paragraph{Stability metrics vs.\ perceptual quality.}
The reported metrics --- boundary cosine similarity, intra-chunk similarity,
identity and object-identity cosines, latent std and band energies ---
measure the \emph{stability} of the latent representation across chunks.
They do not assess perceptual quality, prompt adherence, or realism; the
final outputs of the production runs were verified perceptually by the
authors, and every audio intervention in this paper was gated perceptually
(Section~\ref{sec:exp_scale}).

\paragraph{Hardware specificity.}
The GGUF 4-bit quantisation and CPU block-streaming configuration was
designed for a specific VRAM budget (16\,GB).  Performance, latency, and
numerical behaviour may differ on GPUs with more VRAM (where block-streaming
is unnecessary) or with different quantisation schemes.

%--------------------------------------------------------------------------
\section{Related Work}
\label{sec:related}
%--------------------------------------------------------------------------

\paragraph{Long video generation.}
StreamingT2V~\cite{henschel2024streamingt2v} and FIFO-Diffusion~\cite{kim2024fifo}
generate long videos by sliding-window denoising.  Gen-L-Video~\cite{wang2023gen}
uses temporal attention patching.  FreeNoise~\cite{qiu2024freenoise} reschedules
the initial noise (window-based fusion of a shuffled, reused noise sequence)
to extend a pretrained model's temporal attention window without retraining or
sliding-window denoising at all --- a tuning-free approach to the same
underlying problem (fixed training-time context length) that CLSS addresses by
chunking with closed-loop correction instead.  VidToMe~\cite{li2024vidtome}
merges self-attention tokens across frames to improve temporal consistency in
zero-shot video editing, operating within a single generation pass rather than
across autoregressively-generated chunks.  CLSS differs from all of these by
applying closed-loop corrections specifically designed for the latent-streaming
regime of large audio-video models, where each chunk is a genuinely separate
denoising process conditioned on the previous chunk's output --- the
exposure-bias setting FreeNoise and VidToMe do not target.

\paragraph{Exposure bias in diffusion.}
The exposure bias problem in autoregressive models is well-studied~\cite{bengio2015scheduled}.
For diffusion models, Ning et al.~\cite{ning2023elucidating} and
Salimans \& Ho~\cite{salimans2022progressive} address related distributional
mismatch issues during training, but inference-time corrections for streaming
have received less attention.

\paragraph{Statistical normalisation for consistency.}
Adaptive Instance Normalisation (AdaIN)~\cite{huang2017arbitrary} was introduced
for style transfer, where the goal is to match the mean and standard deviation of
a content feature map to those of a style reference.  In CLSS, AdaIN is
repurposed as a \emph{drift correction} applied to denoised latents between
chunks, using the cross-chunk EMA as the style reference.  The combination of
EMA tracking with AdaIN blending (rather than full replacement) is motivated by
the streaming setting, where the reference must itself evolve to allow
intentional scene changes.

\paragraph{Frequency-domain video editing.}
Yang et al.~\cite{yang2024fresco} and Ceylan et al.~\cite{ceylan2023pix2video}
use frequency-domain constraints for temporal consistency in video editing.
CLSS applies frequency shrinkage in the \emph{latent} domain to a streaming
generation setting, targeting energy accumulation rather than pixel-domain
flicker.

%--------------------------------------------------------------------------
\section{Conclusion}
\label{sec:conclusion}
%--------------------------------------------------------------------------

We presented CLSS, a five-component closed-loop streaming synthesis framework
that enables stable long audio-video generation with LTX-2.3 22B on consumer
16\,GB VRAM hardware.  The corrections --- calibrated context re-noising,
EMA-tracked AdaIN, frequency-band soft shrinkage, a dynamic anchor bank, and a
two-band spatial detail anchor --- are applied between chunks without touching
the transformer weights.  At the production operating point (four 52-lf
chunks, $\approx17.3$\,s windows inside the model's 20-second positional
wall), the pipeline holds video latent std to $\Delta<0.005$, seam boundary
similarity at $0.93$--$0.99$, and the high-frequency energy share flat
end-to-end in both conditioning modes, while the open/closed-loop ablation
shows the corrections are exactly what stands between that stability and
open-loop drift ($+0.118$ std, $+34\%$ audio RMS, $1.60\times$ audio
high-frequency runaway over five chunks).

Two system-level problems were resolved by measured design rather than by
tuning.  The sigma-schedule conflict --- the video requires the scheduler's
$4.68$ shift at 52-lf chunks, the audio the $1.844$ shift --- is resolved by
splitting the schedule per modality inside the single joint pass (shape-only,
full-noise start, auto-derived $0.394$ multiplier), restoring the audio's
latent scale without touching the video's byte-identical sigmas.  The
Stage-2 refinement asymmetry is resolved by the join sigma: audio joins the
video's low-$\sigma$ descent at $\sigma=0.55$, keeping $45\%$ of the
Stage-1 audio and holding its fidelity at $0.66$--$0.74$, where full
re-invention destroys it ($0.07$--$0.16$) because melody and timbre are not
text-specified.  The audio-scale probe (real music unit variance vs.\
generated half-scale, dynR as the one perceptually validated axis) gates every audio
intervention, and the pipeline's per-chunk telemetry --- settings dumps,
per-frame drift-event tables, phase-lock detection, seam-grid object
identity --- localises every failure mode reported here in wall-clock time
and latent-frame position.

Two open problems remain.  First, coarse-layout jump events: the layout
metric and the seam-pin safeguard that target them are in place
(Section~\ref{sec:events}), but their behaviour against a genuine event is
unmeasured, since none occurs in the runs reported here.  Second,
object-level identity persistence: the pooled metrics are object-blind, and
the decoded seam-grid metric observes object swaps without correcting them.
Both are active directions.

%--------------------------------------------------------------------------
\section*{Acknowledgements}
%--------------------------------------------------------------------------
We thank the LTX-Video team at Lightricks for releasing LTX-2.3 and the
associated GGUF weights, and the ComfyUI community for infrastructure.

%--------------------------------------------------------------------------
\bibliographystyle{plain}
\begin{thebibliography}{99}

\bibitem{qiu2024freenoise}
H.~Qiu, M.~Xia, Y.~Zhang, Y.~He, X.~Wang, Y.~Shan, and Z.~Liu.
\newblock {FreeNoise}: Tuning-free longer video diffusion via noise rescheduling.
\newblock In \emph{ICLR}, 2024.

\bibitem{li2024vidtome}
X.~Li, C.~Ma, X.~Yang, and M.-H.~Yang.
\newblock {VidToMe}: Video token merging for zero-shot video editing.
\newblock In \emph{CVPR}, 2024.

\bibitem{henschel2024streamingt2v}
R.~Henschel, L.~Khachatryan, D.~Hayrapetyan, H.~Poghosyan, R.~Abrahamyan,
A.~Nersisyan, S.~I.~Yeranosyan, and O.~Kacher.
\newblock {StreamingT2V}: Consistent, dynamic, and extendable long video generation from text.
\newblock \emph{arXiv:2403.14773}, 2024.

\bibitem{kim2024fifo}
J.~Kim, J.~Jeong, S.~Park, M.~Kim, S.~Kim, and J.~Ye.
\newblock {FIFO-Diffusion}: Generating infinite videos from text without training.
\newblock \emph{arXiv:2405.11473}, 2024.

\bibitem{wang2023gen}
F.~Wang, W.~Chen, G.~Song, H.~Zhang, W.~Li, C.~Shen, and T.~Mei.
\newblock Gen-{L}-Video: Multi-text conditioned long video generation via temporal co-denoising.
\newblock \emph{arXiv:2305.18264}, 2023.

\bibitem{bengio2015scheduled}
S.~Bengio, O.~Vinyals, N.~Jaitly, and N.~Shazeer.
\newblock Scheduled sampling for sequence prediction with recurrent neural networks.
\newblock In \emph{NeurIPS}, 2015.

\bibitem{ning2023elucidating}
M.~Ning, M.~Li, J.~Su, A.~A.~Salah, and I.~Ertugrul.
\newblock Elucidating the exposure bias in diffusion models.
\newblock \emph{arXiv:2308.15321}, 2023.

\bibitem{salimans2022progressive}
T.~Salimans and J.~Ho.
\newblock Progressive distillation for fast sampling of diffusion models.
\newblock In \emph{ICLR}, 2022.

\bibitem{huang2017arbitrary}
X.~Huang and S.~Belongie.
\newblock Arbitrary style transfer in real-time with adaptive instance normalization.
\newblock In \emph{ICCV}, 2017.

\bibitem{yang2024fresco}
S.~Yang, Y.~Zhou, Z.~Liu, and C.~C.~Loy.
\newblock {FRESCO}: Spatial-temporal correspondence for zero-shot video translation.
\newblock \emph{arXiv:2403.12962}, 2024.

\bibitem{ceylan2023pix2video}
D.~Ceylan, C.-H.~P.~Huang, and N.~J.~Mitra.
\newblock Pix2video: Video editing using image diffusion.
\newblock In \emph{ICCV}, 2023.

\end{thebibliography}

\end{document}
